% Complete Section 6.1, original PDF 73-78 / printed 63-68, including all Table 6.1 cells.
% block: exterior-algebra-heading page:73 type:section
\section{{外代数}}\label{sec:exterior-algebra}

% block: exterior-algebra-introduction page:73 type:ordinary-paragraph
{外微积分的主要研究对象是微分 $k$-形式, 下一节将讨论它们. 这些函数的值是交错多重线性形式, 也称为代数 $k$-形式. 外微积分自然涉及这些函数的微分和积分运算. 但首先需要建立每个点处出现的代数 $k$-形式的基本术语和性质, 正如在研究向量微积分之前, 必须先理解有限维向量空间一样. 这正是外代数的内容, 它由 Grassman 在 19 世纪中叶创立. 这里简要回顾其定义和主要性质, 这些都是初等的.}

% block: multilinear-forms-heading page:73 type:subsection
\subsection{{多重线性形式}}\label{subsec:multilinear-forms}

% block: multilinear-forms-definition pages:73,74 type:ordinary-paragraph
{对有限维向量空间 $V$ 和非负整数 $k$, 考虑由所有 $k$-线性映射构成的向量空间 $\operatorname{Lin}^kV$:}
\[
\omega:\overbrace{V\times\cdots\times V}^{k\text{ times}}\to\mathbb R.
\]
{按约定, 将 $\operatorname{Lin}^0V$ 理解为标量空间 $\mathbb R$. 若 $\omega\in\operatorname{Lin}^jV$, $\mu\in\operatorname{Lin}^kV$, 则张量积 $\omega\otimes\mu\in\operatorname{Lin}^{j+k}V$ 由下式给出:}
\[
\begin{gathered}
(\omega\otimes\mu)(v_1,\ldots,v_{j+k})=\omega(v_1,\ldots,v_j)\mu(v_{j+1},\ldots,v_{j+k}),\\
v_1,\ldots,v_{j+k}\in V.
\end{gathered}
\]

% block: multilinear-dual-basis page:74 type:ordinary-paragraph
{由定义, 空间 $\operatorname{Lin}^1V$ 就是对偶空间 $V'$. 其元素称为\emph{余向量}或 $1$-形式. 若选取 $V$ 的一组基 $u_1,\ldots,u_n$ ($n=\dim V$), 就诱导了 $V'$ 的对偶基 $u^1,\ldots,u^n$, 由 $u^i(u_j)=\delta_{ij}$ 定义. $\operatorname{Lin}^kV$ 的一般元素由它赋予每个 $k$ 元组 $(u_{\sigma_1},\ldots,u_{\sigma_k})$ 的值决定, 其中 $\sigma_i$ 是 $1$ 到 $n$ 之间的整数. 因此 $\dim\operatorname{Lin}^kV=n^k$. 将 $k$ 元组 $(u_{\sigma_1},\ldots,u_{\sigma_k})$ 映为 $1$, 并将所有其他 $k$ 元组 $(u_{\tau_1},\ldots,u_{\tau_k})$ 映为 $0$ 的 $\operatorname{Lin}^kV$ 元素, 就是 $u^{\sigma_1}\otimes\cdots\otimes u^{\sigma_k}$, 因而所有这些元素构成 $\operatorname{Lin}^kV$ 的一组基.}

% block: multilinear-canonical-coordinate-basis page:74 type:ordinary-paragraph
{情形 $V=\mathbb R^n$ 与一般情形的区别, 仅在于它具有典范基 $e_1,\ldots,e_n$. 相应对偶基记作 $dx^1,\ldots,dx^n$. 因此, $dx^i$ 就是将 $\mathbb R^n$ 中每个向量映为其第 $i$ 个分量的线性泛函. 随着外微积分的展开, 采用这一特别记号 ($dx^i$ 而不是例如 $e^i$) 的原因将变得清楚.}

% block: alternating-forms-heading page:74 type:subsection
\subsection{{交错多重线性形式}}\label{subsec:alternating-multilinear-forms}

% block: alternating-forms-definition page:74 type:ordinary-paragraph
{交换两个变量时变号的多重线性形式 $\omega\in\operatorname{Lin}^kV$, 称为\emph{反对称}或\emph{交错}形式:}
\[
\begin{gathered}
\omega(v_1,\ldots,v_i,\ldots,v_j,\ldots,v_k)=-\omega(v_1,\ldots,v_j,\ldots,v_i,\ldots,v_k),\\
1\leq i<j\leq n,\qquad v_1,\ldots,v_k\in V.
\end{gathered}
\]
{特别地, 若交错形式的任一自变量重复出现, 则其值为零. 任意多重线性形式 $\omega$ 都伴随一个交错形式, 即其反对称部分:}
\[
(\operatorname{skw}\omega)(v_1,\ldots,v_k)=\frac1{k!}\sum_{\sigma\in\Sigma_k}\operatorname{sign}(\sigma)\omega(v_{\sigma_1},\ldots,v_{\sigma_k}),
\]
{其中对整数 $1$ 到 $k$ 的所有排列求和. 例如, 若 $\omega\in\operatorname{Lin}^3V$, 则}
\[
\begin{aligned}
(\operatorname{skw}\omega)(v_1,v_2,v_3)=\frac16[&\omega(v_1,v_2,v_3)+\omega(v_2,v_3,v_1)+\omega(v_3,v_1,v_2)\\
&-\omega(v_1,v_3,v_2)-\omega(v_3,v_2,v_1)-\omega(v_2,v_1,v_3)].
\end{aligned}
\]

% block: alternating-space-dimension page:74 type:ordinary-paragraph
{将空间 $\operatorname{Alt}^kV$ 定义为所有反对称 $k$-线性形式构成的空间. 空间 $\operatorname{Alt}^0V$ 和 $\operatorname{Alt}^1V$ 分别与 $\operatorname{Lin}^0V$ 和 $\operatorname{Lin}^1V$ 一致, 即分别是标量空间和余向量空间. $\operatorname{Alt}^kV$ 的一个元素, 由它赋予满足 $1\leq\sigma_1<\cdots<\sigma_k\leq n$ 的每个 $k$ 元组 $(e_{\sigma_1},\ldots,e_{\sigma_k})$ 的值决定, 所以 $\dim\operatorname{Alt}^kV=\binom nk$. 当 $k>n$ 时, $\operatorname{Alt}^kV=0$.}

% block: exterior-product-definition pages:74,75 type:ordinary-paragraph
{$\omega\in\operatorname{Alt}^jV$ 和 $\mu\in\operatorname{Alt}^kV$ 的\emph{外积}或\emph{楔积}由下式给出:}
\[
(\omega\wedge\mu)(v_1,\ldots,v_{j+k})=\sum_\sigma\operatorname{sign}(\sigma)\omega(v_{\sigma_1},\ldots,v_{\sigma_j})\mu(v_{\sigma_{j+1}},\ldots,v_{\sigma_{j+k}}),
\]
{其中对满足 $\sigma_1<\cdots<\sigma_j$ 和 $\sigma_{j+1}<\cdots<\sigma_{j+k}$ 的 $\sigma\in\Sigma_{j+k}$ 求和. 等价地,}
\[
\omega\wedge\mu=\binom{j+k}j\operatorname{skw}(\omega\otimes\mu).
\]
{外积满足反交换律}
\[
\omega\wedge\mu=(-)^{jk}\mu\wedge\omega,\qquad\omega\in\operatorname{Alt}^jV,\ \mu\in\operatorname{Alt}^kV.
\]

% block: alternating-basis page:75 type:ordinary-paragraph
{设 $\{u_i\}$ 是 $V$ 的一组基, $\{u^i\}$ 是相应的对偶基. 给定 $1\leq\sigma_1<\cdots<\sigma_k\leq n$, 将 $k$ 元组 $(u_{\sigma_1},\ldots,u_{\sigma_k})$ 映为 $1$, 将所有其他具有递增指标的 $k$ 元组映为 $0$ 的 $\operatorname{Alt}^kV$ 元素为}
\begin{equation}\label{eq:alternating-form-basis}
u^{\sigma_1}\wedge\cdots\wedge u^{\sigma_k}.
\end{equation}
{所有这样的元素构成 $\operatorname{Alt}^kV$ 的一组基, 因而 $V$ 上的一般交错 $k$-形式是这些元素的线性组合. 特别地, $\mathbb R^n$ 上的一般交错 $k$-形式可写为}
\[
\sum_{1\leq\sigma_1<\cdots<\sigma_k\leq n}a_\sigma\,dx^{\sigma_1}\wedge\cdots\wedge dx^{\sigma_k}
\]
{其中系数 $a_\sigma\in\mathbb R$ 唯一确定.}

% block: alternating-operations-heading page:75 type:subsection
\subsection{{交错形式上的运算}}\label{subsec:alternating-form-operations}

% block: alternating-contraction page:75 type:ordinary-paragraph
\medskip\noindent\textbf{{缩并}}\par
{除 $\operatorname{Alt}^kV$ 上的向量空间运算和外积外, 还有其他几种有用的代数运算. 给定 $\omega\in\operatorname{Alt}^kV$, 其中 $k\geq1$, 以及 $v\in V$, $\omega$ 与 $v$ 的\emph{缩并}, 也称为内乘, 是如下 $(k-1)$-形式:}
\[
(\omega\lrcorner v)(v_1,\ldots,v_{k-1})=\omega(v,v_1,\ldots,v_{k-1}).
\]
{(若 $k=0$, 规定 $\omega\lrcorner v:=0$.) 由于 $\omega$ 是交错的, $(\omega\lrcorner v)\lrcorner v$ 总为零. 此外, 容易验证, 与一个向量的缩并是一种反导子, 即}
\[
\begin{gathered}
(\omega\wedge\eta)\lrcorner v=(\omega\lrcorner v)\wedge\eta+(-1)^k\omega\wedge(\eta\lrcorner v),\\
\omega\in\operatorname{Alt}^kV,\ \eta\in\operatorname{Alt}^lV.
\end{gathered}
\]

% block: alternating-inner-product page:75 type:ordinary-paragraph
\medskip\noindent\textbf{{内积}}\par
{若向量空间 $V$ 赋有\emph{内积} $\langle\cdot,\cdot\rangle$, 就诱导了从 $V$ 到 $V'$ 的 Riesz 同构, 从而在 $V'$ 上诱导内积. 于是可在 $\operatorname{Alt}^kV$ 上定义\emph{内积}}
\[
\langle u^1\wedge\cdots\wedge u^k,v^1\wedge\cdots\wedge v^k\rangle=\det\big[(\langle u^i,v^j\rangle)_{i,j=1}^n\big],
\]
{并通过双线性性延拓到 $\operatorname{Alt}^kV$ 的一般元素. 若 $u_i$ 构成 $V$ 的\emph{标准正交}基, 则相应对偶基元素 $u^i$ 在 $V'$ 中标准正交, 由此得到的 $\operatorname{Alt}^kV$ 的基函数 \eqref{eq:alternating-form-basis} 也标准正交.}

% block: volume-form page:76 type:ordinary-paragraph
\medskip\noindent\textbf{{体积形式}}\par
{当 $k$ 等于 $V$ 的维数 $n$ 时, 空间 $\operatorname{Alt}^nV$ 是一维的, 因而一组基只包含一个元素. 若空间 $V$ 既赋有内积又赋有\emph{定向}, 就有典范的基元素选取. 有限维向量空间 $V$ 的有序基分为两类, 取决于连接其中两组基的换基矩阵的行列式为正还是为负. 指定其中一类的有序基为正向基, 就给向量空间定了向. (在 $\mathbb R^3$ 中, 习惯上将满足"右手法则"的基取为正向基.) 若 $\omega\in\operatorname{Alt}^nV$, 则对任意标准正交正向基 $u_1,\ldots,u_n$, $\omega(u_1,\ldots,u_n)$ 的值相同; 要求此值为 $1$, 就唯一确定了 $\operatorname{Alt}^nV$ 的一个元素. 得到的交错 $n$-形式称为 $V$ 上的\emph{体积形式}, 通常记为 $\operatorname{vol}$. 它可用任意正向标准正交基的对偶基写成 $\operatorname{vol}=u^1\wedge\cdots\wedge u^n$. 在 $\mathbb R^n$ 上, 典范基 $e_1,\ldots,e_n$ 标准正交且正向, 所以}
\[
\operatorname{vol}=dx^1\wedge\cdots\wedge dx^n.
\]

% block: hodge-star-definition page:76 type:ordinary-paragraph
\medskip\noindent\textbf{Hodge {星算子}}\par
{继续考虑有定向的内积空间 $V$. 设 $\omega\in\operatorname{Alt}^kV$, 其中 $0\leq k\leq n$, 并设 $\mu\in\operatorname{Alt}^{n-k}V$. 则对某个实数 $c$, 有 $\omega\wedge\mu=c\operatorname{vol}$, 且映射 $\mu\mapsto c$ 是 $\operatorname{Alt}^{n-k}$ 上的线性泛函. 由 Riesz 表示定理, 它可表示为 $\mu$ 与 $\operatorname{Alt}^{n-k}V$ 中某个其他元素的内积. 该元素依赖于 $\omega$, 称为 $\omega$ 的 \emph{Hodge 星算子}或 Hodge 对偶, 记为 $\star\omega$. 因此 Hodge 星算子由下述关系刻画:}
\begin{equation}\label{eq:hodge-star-characterization}
\omega\wedge\mu=\langle\star\omega,\mu\rangle\operatorname{vol},\qquad\omega\in\operatorname{Alt}^kV,\ \mu\in\operatorname{Alt}^{n-k}V.
\end{equation}

% block: hodge-star-basis-and-square page:76 type:ordinary-paragraph
{因此, 对每个 $k$, 都有线性映射 $\star:\operatorname{Alt}^kV\to\operatorname{Alt}^{n-k}V$. 若选取 $1$-形式的正向标准正交基 $u^1,\ldots,u^n$, 使 $\operatorname{vol}=u^1\wedge\cdots\wedge u^n$, 则有}
\[
\star u^{\sigma_1}\wedge\cdots\wedge u^{\sigma_k}=\pm u^{\tau_1}\wedge\cdots\wedge u^{\tau_{n-k}}
\]
{其中 $\sigma$ 和 $\tau$ 是将 $\{1,\ldots,n\}$ 分割的递增序列, 符号是排列 $(\sigma,\tau)$ 的符号. 例如, 在 $\mathbb R^3$ 中,}
\[
\star dx^1=dx^2\wedge dx^3,\qquad\star dx^2=-dx^1\wedge dx^3,\qquad\star dx^3=dx^1\wedge dx^2.
\]
{注意, $\star$ 是\emph{等距映射}, 且 $\star\star\omega=(-)^{k(n-k)}\omega$; 或等价地, 若 $\omega\in\operatorname{Alt}^kV$, 则}
\[
\star\star\omega=\omega\text{ in odd dimensions},\qquad\star\star\omega=(-)^k\omega\text{ in even dimensions}.
\]
{结合 \eqref{eq:hodge-star-characterization}, 得到 Hodge 星算子的另一种刻画:}
\begin{equation}\label{eq:hodge-star-inner-product}
\omega\wedge\star\nu=\langle\omega,\nu\rangle\operatorname{vol},\qquad\omega,\nu\in\operatorname{Alt}^kV.
\end{equation}

% block: scalar-vector-proxies pages:76,77 type:ordinary-paragraph
\medskip\noindent\textbf{{标量代理和向量代理}}\par
{若 $V$ 是维数为 $n$ 的有定向内积空间, 则空间 $\operatorname{Alt}^0V$ 和 $\operatorname{Alt}^nV$ 各自与标量空间 $\mathbb R$ 之间存在自然同构, 空间 $\operatorname{Alt}^1V$ 和 $\operatorname{Alt}^{n-1}V$ 各自与向量空间 $V$ 之间也存在自然同构. 由此, 当 $k=0$, $1$, $n-1$ 或 $n$ 时 (因而在 $\dim V\leq3$ 时涵盖所有情形), 可为 $\operatorname{Alt}^kV$ 的每个元素关联一个\emph{标量代理}或\emph{向量代理}. 现在描述这些同构.}

% block: proxy-isomorphisms page:77 type:ordinary-paragraph
{对任意向量空间 $V$, $\operatorname{Alt}^0V$ 就是标量空间 $\mathbb R$, 相应同构为恒等映射. $V$ 上的内积通过 Riesz 表示定义了从 $\operatorname{Alt}^1V=V'$ 到 $V$ 上的同构. 随后, Hodge 星算子将 $\operatorname{Alt}^{n-1}V$ 同构地映到 $\operatorname{Alt}^1V$ (后者与 $V$ 同构), 并类似地将 $\operatorname{Alt}^nV$ 映到 $\operatorname{Alt}^0V=\mathbb R$.}

% block: two-dimensional-proxies page:77 type:ordinary-paragraph
{情形 $n=2$ 比较特殊, 因为此时有两种方式将向量代理与交错 $1$-形式关联: 或者将它视为 $1$-形式并应用 Riesz 映射, 或者视为 $(n-1)$-形式, 即先应用 Hodge 星算子再应用 Riesz 映射. 这两种方式并不一致, 因为 Hodge 星算子 $\operatorname{Alt}^1V\to\operatorname{Alt}^1V$ 不是恒等映射: 它将 $u^1$ 映为 $u^2$, 将 $u^2$ 映为 $-u^1$. 因此, 两个向量代理相差一个直角旋转.}

% block: linear-pullback page:77 type:ordinary-paragraph
\medskip\noindent\textbf{{线性映射下的拉回}}\par
{若 $L:V\to W$ 是有限维向量空间之间的线性映射, 则有一个显然的\emph{拉回}运算 $L^*:\operatorname{Alt}^kW\to\operatorname{Alt}^kV$, 定义为}
\begin{equation}\label{eq:linear-pullback-definition}
L^*\omega(v_1,\ldots,v_k)=\omega(Lv_1,\ldots,Lv_k).
\end{equation}
{拉回以反变方式作用: 若 $K:W\to X$ 是另一个映射, 则 $(K\circ L)^*=L^*\circ K^*$. 它也保持代数结构, 包括外积:}
\begin{equation}\label{eq:linear-pullback-wedge}
L^*(\omega\wedge\mu)=L^*\omega\wedge L^*\mu.
\end{equation}

% block: pullback-inclusion-and-normal page:77 type:ordinary-paragraph
{当 $V$ 是 $W$ 的子空间且 $L$ 是包含映射时, $W$ 上一个 $k$-形式的拉回就是它在 $V\times\cdots\times V$ 上的限制. 假设 $W$ 具有内积, 则可通过向量代理将拉回 $\operatorname{Alt}^1W\to\operatorname{Alt}^1V$ 看作从 $W\to V$ 的映射, 它正是正交投影. 另一个例子是, $W$ 为维数 $n$ 的有定向内积空间, $V$ 为维数 $n-1$ 的有定向子空间. 此时可使用代理, 将拉回 $\operatorname{Alt}^{n-1}W\to\operatorname{Alt}^{n-1}V$ 看作从 $W\to\mathbb R$ 的映射. 来计算此映射. 先取 $V$ 的正向标准正交基 $u_2,\ldots,u_n$, 在前面添上 $V$ 在 $W$ 中的单位法向量 $u_1$, 并选取其符号使 $u_1,\ldots,u_n$ 为正向. 用 $u^1,\ldots,u^n$ 表示 $W'$ 的对偶基. $u^1$ 到 $V'$ 的拉回为零, 其余 $u^i$ 的拉回 $v^2,\ldots,v^n$ 则构成 $V'$ 中与 $u_2,\ldots,u_n$ 对偶的基. 由于拉回保持外积, 有}
\[
\begin{gathered}
L^*(u^2\wedge\cdots\wedge u^n)=v^2\wedge\cdots\wedge v^n=\operatorname{vol}_V,\\
L^*(u^1\wedge u^2\wedge\cdots\wedge\widehat{u^i}\wedge\cdots\wedge u^n)=0.
\end{gathered}
\]
{转用代理表示, 拉回映射 $L^*$ 因而成为}
\[
u_1\mapsto1,\qquad u_i\mapsto0,\qquad i=2,\ldots,n.
\]
{此映射就是 $W$ 中向量与有定向单位法向量 $u_1$ 的内积. 总之, 用代理表示时,}

\begin{itemize}
\item{一个 $1$-形式到子空间的拉回就是正交投影;}
\item{一个 $n-1$ 形式到余维为 $1$ 的子空间的拉回, 就是与法向量的内积.}
\end{itemize}

% block: exterior-algebra-r3-table page:78 type:table
\begin{table}[htbp]\centering\footnotesize
\caption{{用标量和向量代理场表示的 $\mathbb R^3$ 外代数.}}\label{tab:exterior-algebra-r3}
\setlength{\tabcolsep}{4pt}
\begin{tabular}{l|l}\hline
\multicolumn{2}{c}{{代理}}\\\hline
$\operatorname{Alt}^0\mathbb R^3=\mathbb R$&$c\leftrightarrow c$\\
$\operatorname{Alt}^1\mathbb R^3\xrightarrow{\cong}\mathbb R^3$&$u_1dx^1+u_2dx^2+u_3dx^3\leftrightarrow u$\\
$\operatorname{Alt}^2\mathbb R^3\xrightarrow{\cong}\mathbb R^3$&$u_1dx^2\wedge dx^3-u_2dx^1\wedge dx^3+u_3dx^1\wedge dx^2\leftrightarrow u$\\
$\operatorname{Alt}^3\mathbb R^3\xrightarrow{\cong}\mathbb R$&$c\leftrightarrow c\,dx^1\wedge dx^2\wedge dx^3$\\\hline
\multicolumn{2}{c}{{外积}}\\\hline
$\wedge:\operatorname{Alt}^1\mathbb R^3\times\operatorname{Alt}^1\mathbb R^3\to\operatorname{Alt}^2\mathbb R^3$&$\times:\mathbb R^3\times\mathbb R^3\to\mathbb R^3$\\
$\wedge:\operatorname{Alt}^1\mathbb R^3\times\operatorname{Alt}^2\mathbb R^3\to\operatorname{Alt}^3\mathbb R^3$&$\cdot:\mathbb R^3\times\mathbb R^3\to\mathbb R$\\\hline
\multicolumn{2}{c}{{与向量 $v\in\mathbb R^3$ 的缩并}}\\\hline
$\lrcorner v:\operatorname{Alt}^1\mathbb R^3\to\operatorname{Alt}^0\mathbb R^3$&$v\cdot:\mathbb R^3\to\mathbb R$\\
$\lrcorner v:\operatorname{Alt}^2\mathbb R^3\to\operatorname{Alt}^1\mathbb R^3$&$v\times:\mathbb R^3\to\mathbb R^3$\\
$\lrcorner v:\operatorname{Alt}^3\mathbb R^3\to\operatorname{Alt}^2\mathbb R^3$&$v:\mathbb R\to\mathbb R^3$\quad$(c\mapsto cv)$\\\hline
\multicolumn{2}{c}{{内积}}\\\hline
{$\operatorname{Alt}^k\mathbb R^3$ 上的内积}&{$\mathbb R$ 或 $\mathbb R^3$ 上的 Euclidean 内积}\\\hline
\multicolumn{2}{c}{{体积形式}}\\\hline
$\operatorname{vol}=dx^1\wedge dx^2\wedge dx^3$&$(v_1,v_2,v_3)\mapsto\det[v_1|v_2|v_3]$\\\hline
\multicolumn{2}{c}{Hodge {星算子}}\\\hline
$\star:\operatorname{Alt}^0\mathbb R^3\to\operatorname{Alt}^3\mathbb R^3$&$\operatorname{id}:\mathbb R\to\mathbb R$\\
$\star:\operatorname{Alt}^1\mathbb R^3\to\operatorname{Alt}^2\mathbb R^3$&$\operatorname{id}:\mathbb R^3\to\mathbb R^3$\\
$\star:\operatorname{Alt}^2\mathbb R^3\to\operatorname{Alt}^1\mathbb R^3$&$\operatorname{id}:\mathbb R^3\to\mathbb R^3$\\
$\star:\operatorname{Alt}^3\mathbb R^3\to\operatorname{Alt}^0\mathbb R^3$&$\operatorname{id}:\mathbb R\to\mathbb R$\\\hline
\multicolumn{2}{c}{{线性映射 $L:\mathbb R^3\to\mathbb R^3$ 的拉回}}\\\hline
$L^*:\operatorname{Alt}^0\mathbb R^3\to\operatorname{Alt}^0\mathbb R^3$&$\operatorname{id}:\mathbb R\to\mathbb R$\\
$L^*:\operatorname{Alt}^1\mathbb R^3\to\operatorname{Alt}^1\mathbb R^3$&$L^{\mathrm T}:\mathbb R^3\to\mathbb R^3$\\
$L^*:\operatorname{Alt}^2\mathbb R^3\to\operatorname{Alt}^2\mathbb R^3$&$(\det L)L^{-1}:\mathbb R^3\to\mathbb R^3$\\
$L^*:\operatorname{Alt}^3\mathbb R^3\to\operatorname{Alt}^3\mathbb R^3$&$(\det L):\mathbb R\to\mathbb R$\quad$(c\mapsto c\det L)$\\\hline
\end{tabular}
\end{table}

% block: exterior-algebra-r3-summary page:78 type:ordinary-paragraph
\medskip\noindent\textbf{{$\mathbb R^3$ 的外代数}}\par
{由于空间 $\operatorname{Alt}^0\mathbb R^n$ 和 $\operatorname{Alt}^n\mathbb R^n$ 可与 $\mathbb R$ 等同, 空间 $\operatorname{Alt}^1\mathbb R^n$ 和 $\operatorname{Alt}^{n-1}\mathbb R^n$ 可与 $\mathbb R^n$ 等同, 在 $n=3$ (或更低) 维时, 外代数的所有运算都可视为标量或向量上的运算. Table~\ref{tab:exterior-algebra-r3} 列出了这些运算, 该表改编自 \BookCite{11}. 注意, 表中倒数第二行作为 $2$-形式拉回出现的矩阵 $(\det L)L^{-1}$ 是 $L$ 的伴随矩阵, 即使 $L$ 奇异, 它也可定义为代数余子式矩阵的转置.}
